\chapter{Introdução}

A introdução é a seção inicial do trabalho acadêmico e tem como objetivo apresentar de forma clara o tema, o contexto e a importância do projeto. Ela prepara o leitor para compreender o desenvolvimento do estudo.

\section{Contexto}
O contexto fornece uma breve descrição do cenário em que o projeto está inserido.  
Aqui você explica o ambiente, os problemas existentes, tendências tecnológicas ou situações práticas que justificam a necessidade do projeto.  
\textbf{Exemplo:} "O aumento do uso de aplicativos móveis exige soluções eficientes para gerenciamento de tarefas em dispositivos Android e iOS."  Aqui você insere suas referências!
\textbf{É aqui que você demonstra seu embasamento teórico e conhecimento sobre  tema usando referências fortes (artigos científicos ou livros)!
}\section{Justificativa}
A justificativa explica por que o projeto é relevante para o campo da engenharia de software.  
Ela responde perguntas como: “Por que este projeto merece ser desenvolvido?” ou “Que problema real ele resolve?”  
\textbf{Exemplo:} "Desenvolver um sistema de gerenciamento de tarefas integrado pode melhorar a produtividade de equipes de desenvolvimento e contribui para melhores práticas de engenharia de software."

\section{Objetivos}
Os objetivos definem o que o projeto pretende alcançar.  
Podem ser divididos em:

\begin{itemize}
    \item \textbf{Objetivo Principal:} 
    % Descreva aqui a meta central do seu projeto
    Ex.: "Desenvolver um aplicativo móvel multiplataforma para gerenciamento de tarefas."

    \item \textbf{Objetivos Secundários:} 
    % Liste aqui os objetivos complementares que ajudam a atingir o objetivo principal
    Ex.: 
    \begin{itemize}
        \item Definir o modelo de banco de dados
        \item Desenhar os protótipos de telas
        \item Configurar ambiente de versionamento
        \item Implementar telas e interação com usuário
        \item Implementar persistência de dados
        \item Configurar ambiente de implantação
        % Adicione outros objetivos secundários conforme necessário
    \end{itemize}
\end{itemize}

\textbf{Resumo:} A introdução deve contextualizar o projeto, justificar sua relevância e definir claramente os objetivos a serem alcançados.






